%%%%%%%%%%%
\section{Graded unitarity structures on affine Kac--Moody algebras} \label{sec:grunitstructures}
%%%%%%%%%%%
Any vertex algebra associated to a unitary four-dimensional $\mathcal{N}=2$ SCFT comes equipped with a graded unitary structure. This is a package of extra structures that consists of a (generally order-four) anti-linear conjugation $\rho$ as well as a (generally half-integral) nondegenerate \emph{good} filtration $\mfR_\bullet$. The $\mfR$-filtrations on the VOAs associated with interacting four-dimensional SCFTs must obey additional conditions, and were dubbed \textit{interacting $\mfR$-filtrations (of four-dimensional type)} in \cite{ArabiArdehali:2025fad}. 

We are interested in cases where the underlying vertex algebra is taken to be the affine Kac--Moody algebra $V_k(\gf)$, where $\gf$ is a semisimple Lie algebra. While the bulk of this paper is dedicated to the case where $\gf$ is simply laced, we make no such restriction here, and discuss the graded unitarity structures more generally. Since $V_k(\gf)$ is a $\mathbb{Z}$-graded bosonic vertex algebra, the graded-unitary version of the spin-statistics relation \cite{Beem:2025guj} requires that the $\mfR$-filtration be integral. As a consequence, the conjugation $\rho$ must be an involution.

Since $\rho$ is a $\mathbb{C}$-antilinear vertex algebra automorphism, its action on the vertex algebra is completely determined by its action on the strong generators $J^a$. Furthermore, the action of $\rho$ selects the compact real form, and with a suitable choice of basis, its action on the strong generators can be expressed as $\rho(J^a) = -J^a$. The determination of the full $\mathfrak{R}$-filtration on $V_k(\gf)$ from first principles is beyond the scope of this work. Instead, we \emph{assume} that the $\mathfrak{R}$-filtration is the one proposed in \cite{ArabiArdehali:2025fad},
\begin{equation} \label{eq:AKMRfiltration}
\scalebox{0.9}{$\displaystyle
    \mathfrak{R}_p V_k(\gf) = {\rm span}\Big\{ (\prod_{i=1}^{\text{rank}~\gf} C^{(s_i)}_{-l_1^{(i)}} \cdots C^{(s_i)}_{-l_{k_i}^{(i)}}) J^{a_1}_{-n_1} \cdots J^{a_k}_{-n_k}  \Omega~; l_j^{(i)} \geqslant s_i, n_j \geqslant 1; ~~k+\sum_{i=1}^{\text{rank}~\gf}(s_i-1) k_i \leqslant p \Big\}~,
    $}
\end{equation}
where $C^{(s_i)}$ are the Casimir operators of \cite{Bouwknegt:1992wg},\footnote{Throughout the paper, the normal ordering is taken to be the nested normal ordering.}
\begin{equation}
    C^{(s_i)}(z) = d_{a_1 a_2 \ldots a_{s_i}} (J^{a_1} J^{a_2} \cdots J^{a_{s_i}})(z), \qquad i = 1, 2, \ldots, \text{rank}~\gf~,
\end{equation}
where $s_i$ are the degrees of primitive Casimir invariants of $\gf$ and $d_{a_1a_2\ldots a_{s_i}}$ are primitive symmetric invariant tensors in $\gf$. Assuming that the $\mfR$-filtration is non-degenerate, one obtains the orthogonal decomposition of a fixed conformal weight subspace, 
\begin{equation} \label{eq:orthogdecomp}
    \left(V_k(\gf) \right)_h = \bigoplus_{R = 0}^h \left(V_k(\gf) \right)_{h,R}
\end{equation}
with the upper bound following from the Schur condition $h = R+j_1+j_2$, which implies $R \leqslant h$.

The $\mathfrak{R}$-filtration in \eqref{eq:AKMRfiltration} is then the one that assigns each Casimir operator $C^{(s_i)}$ to the filtration component $p = s_i - 1$. This $\mathfrak{R}$-filtration is the ``simplest" modification consistent with the Higgs branch relations coming from the 4d $\mathcal{N} = 2$ SCFT which have $V_k(\gf)$ as the associated VOAs. This standard $\mfR$-filtration is also the one that was implicitly used in the computations of Macdonald indices \cite{Gadde:2011uv} carried out in \cite{Song:2016yfd, kang2026indexnonpointrgflows, Kang:2026ney}.\footnote{See \cite{Jiang:2026ffv} for a discussion on using the $\mfR$-filtration and graded unitarity structures for computing specific limits of Macdonald indices.}  

%%%%%%%%%%
\subsection{Goodness of the standard $\mathfrak{R}$-filtration}
%%%%%%%%%%
The $\mathfrak{R}$-filtration on any vertex algebra associated to a 4d $\mathcal{N} = 2$ SCFT must satisfy the necessary property of \emph{goodness}, which follows from noting that the OPE of two Schur operators is non-singular. Additionally, the $\mfR$-filtration is increasing and exhaustive. Concretely, for affine Kac--Moody algebras, an integral, increasing, and exhaustive good filtration is \cite{LiVertexandVPA, ArabiArdehali:2025fad},
\begin{equation}
    \ldots \subseteq \mathcal{G}_{-1} \mathbb{V} \subseteq \mathcal{G}_{0} \mathbb{V} \subseteq \mathcal{G}_{1} \mathbb{V} \subseteq \ldots, \qquad \bigcup_{m \in \mathbb{Z}}^\infty \mathcal{G}_{m} \mathbb{V} = \mathbb{V}~,
\end{equation}
for which the identity $\Omega \in \mathcal{G}_0 \mathbb{V}$, and for any $a \in \mathcal{G}_{m}\mathbb{V}$ and $b \in \mathcal{G}_{n}\mathbb{V}$,
\begin{equation}
\begin{split}
    a_i b \in \mathcal{G}_{m+n} \mathbb{V}, \qquad i < 0~, \\
    a_i b \in \mathcal{G}_{m+n-1} \mathbb{V}, \ \ \  i \geqslant 0~.
\end{split}
\end{equation}
Taking the commutative limit (more precisely, the associated graded) of the VOA with respect to any good filtration yields a Poisson vertex algebra (PVA) \cite{LiVertexandVPA}. The PVA obtained by taking the commutative limit with respect to the $\mfR$-filtration corresponds to the PVA obtained by the holomorphic-topological twist of the parent 4d theory \cite{Kapustin:2006hi}. Conversely, the VOA can be viewed as a certain Omega deformation of the holomorphically twisted theory \cite{Oh:2019bgz, Jeong:2019pzg}. 

It was conjectured in \cite{ArabiArdehali:2025fad} that the standard $\mathfrak{R}$-filtration in \eqref{eq:AKMRfiltration} is \emph{good}. In the remainder of this section, we will show that this is indeed the case. It suffices to show that all operators appearing in the singular terms of the OPEs $J^a(z) C^{(s_i)}(0)$ and $C^{(s_i)}(z) C^{(s_j)}(0)$ are in the filtered subspaces  $\mathfrak{R}_{s_i-1} V_k(\gf)$ and $\mathfrak{R}_{s_i+s_j-3} V_k(\gf)$ respectively. To this end, note that the relevant OPEs have the following schematic form,
\begin{equation}
\begin{split}
    J^a(z) C^{(s_i)}(0) = \sum_{k = 1}^{s_i+1} \frac{\mathcal{A}_k}{z^k} + \text{regular terms}~, \\
    C^{(s_i)}(z) C^{(s_j)}(0) = \sum_{k = 1}^{s_i+s_j} \frac{\mathcal{B}_k}{z^k} + \text{regular terms}~.
\end{split}
\end{equation}
Here, $\mathcal{A}_k, \mathcal{B}_k$ are sums of derivatives of normal ordered products of the currents at the relevant pole order. Conformal covariance requires that $h(\mathcal{A}_k) = s_i-k+1 \leqslant s_i -1$ for all $k \geqslant 2$, and $h(\mathcal{B}_k) = s_i+s_j-k \leqslant s_i + s_j - 3$ for all $k \geqslant 3$. As a result, it follows from \eqref{eq:orthogdecomp} that $\mathcal{A}_{k \geqslant 2} \in \mathfrak{R}_{s_i - 1}V_k(\gf)$, and $\mathcal{B}_{k \geqslant 3} \in \mathfrak{R}_{s_i + s_j -3}V_k(\gf)$. Since the Casimir operators are constructed using invariant tensors, the order one pole term in the $J^aC^{(s_i)}$ OPE vanishes, implying $\mathcal{A}_{k \geqslant 1} \in \mathfrak{R}_{s_i - 1}V_k(\gf)$. Thus, proving goodness of the filtration reduces to showing that the operators appearing in the $C^{(s_i)} C^{(s_j)}$ OPE at the order one pole with \emph{zero} or \emph{one} derivatives, and order two pole with \emph{zero} derivatives belong to the filtered subspace $\mathfrak{R}_{s_i + s_j -3}V_k(\gf)$. 

It is convenient to consider another filtration that keeps track of derivatives: the Li filtration. All vertex algebras are equipped with the Li filtration, which is a decreasing, good filtration. For affine Kac--Moody algebras $V_k(\gf)$, this filtration is given by,
\begin{equation}
    \mathfrak{F}^p V_k(\gf)= \mathrm{span}\left\{ J^{a_1}_{-1-n_1} J^{a_2}_{-1-n_2} \ldots J^{a_m}_{-1-n_m} \Omega; \ \ \ n_i \geqslant 0; \ \ \ \sum_{i = 1}^{m} n_i \geqslant p \right\}~.
\end{equation}
The degree-zero subspace in the associated graded of $V_k(\gf)$ with respect to the Li filtration comprises sums of normal ordered products of currents without any derivatives, $\mathrm{gr}_0^{\mathfrak{F}^\bullet} V_k(\gf) \cong \mathfrak{F}^0V_k(\gf) / \mathfrak{F}^1V_k(\gf)$, and is also referred to as the Zhu $C_2$ algebra \cite{Zhu:1996gaq}. Similarly, the quotient algebra $\mathfrak{F}^1V_k(\gf) / \mathfrak{F}^2V_k(\gf)$ comprises sums of normal ordered products of operators with exactly one derivative. The strategy to proving the goodness of the $\mfR$-filtration in \eqref{eq:AKMRfiltration} is to use the Li filtration and reduce the problem to a question of invariant Lie theory. 

To begin with, consider the OPE of the currents to be given by,
\begin{equation}
    J^a(z) J^b(w) = \frac{k \delta^{a b}}{(z-w)^2} + \frac{i f^{ab}_{\ \ c} J^c(w)}{(z-w)} + \text{regular terms}~,
\end{equation}
where we have chosen the normalisation such that the non-degenerate Killing form on the Lie algebra $\gf$ is $\kappa^{ab} = \delta^{ab}$. We will also use the resulting identity,
\begin{equation} \label{eq:AntiSymJJ}
    (J^a J^b) - (J^bJ^a) = if^{ab}_{\ \ c} \partial J^c~. 
\end{equation}
Following standard techniques from \cite{Bais:1987dc, SEVRIN1988465, Kac:1996wd}, it can be shown by direct computation that the $\mathcal{O}(z^{-2})$ pole term is given by,
\begin{equation} \label{eq:order2pole}
\begin{split}
    C^{(s_i)}(z) C^{(s_j)}(0)|_{\mathcal{O}(z^{-2})} &= d_{a_1 \ldots a_{s_i}} d_{b_1 \ldots b_{s_j}}\left(k \sum_{m=1}^{s_i} \sum_{n=1}^{s_j} \delta^{a_m,b_n} \left(J^{\hat{A}_m} J^{\hat{B}_n}\right) \right. \\ 
    &\left. - \sum_{\substack{m=1}}^{s_i} \sum_{1 \leqslant n_1 < n_2}^{s_j} f^{a_{m} b_{n_1}}_{\qquad \ c} f^{c b_{n_2}}_{\ \ \ \ \ d} \left(J^d J^{\hat{A}_{m}} J^{\hat{B}_{(n_1,n_2)}}\right) \right) + \ldots~, 
\end{split}
\end{equation}
where the shorthand notations indicate,
\begin{equation}
\begin{split}
    &J^{\hat{A}_{m}} := J^{a_1} \cdots J^{a_{m-1}} J^{a_{m+1}} \cdots J^{a_{s_i}}, \qquad J^{\hat{B}_{n}} := J^{b_1} \cdots J^{b_{n-1}} J^{b_{n+1}} \cdots J^{b_{s_j}}~, \\
    &J^{\hat{B}_{(n_1,n_2)}} := J^{b_1} \cdots J^{b_{n_1-1}} J^{b_{n_1+1}} \cdots J^{b_{n_2-1}} J^{b_{n_2+1}} \cdots J^{b_{s_j}}~,
\end{split}
\end{equation}
and $\ldots$ represent sums of normal ordered products of operators with at least one derivative, and hence belong to $\mathfrak{F}^1 V_k(\gf)$. Since re-normal ordering of the operators introduces at least one more derivative, the tensor appearing in the coefficient of the normal ordered products can be completely symmetrised in all the indices without changing the $\mathcal{O}(z^{-2})$ pole term in the Zhu $C_2$ algebra. Therefore, the different terms can all be clubbed together and rewritten in the schematic form,
\begin{equation} \label{eq:order2simp}
    C^{(s_i)}(z) C^{(s_j)}(0)|_{\mathcal{O}(z^{-2})} = X^{(1)}_{(\alpha_1 \alpha_2 \ldots \alpha_{s_i+s_j-2})} (J^{\alpha_1} J^{\alpha_2} \ldots J^{\alpha_{s_i+s_j-2}}) + \ldots~,
\end{equation}
where $\ldots$ represent sums of normal ordered products of operators with at least one derivative, and hence belong to $\mathfrak{F}^1 V_k(\gf)$. The tensor $X^{(1)}_{(\alpha_1 \alpha_2 \ldots \alpha_{s_i+s_j-2})}$ is a fully symmetric invariant tensor, with the invariance following from the $G$-equivariance of the OPE, and is hence valued in the symmetric algebra $\left(S^{(s_i+s_j-2)} \gf^*\right)^{\gf}$. Therefore, the $\mathcal{O}(z^{-2})$ term in the OPE with no derivatives can be expressed as sums of normal ordered products of Casimir operators, leading to,
\begin{equation}
    C^{(s_i)}(z) C^{(s_j)}(0)|_{\mathcal{O}(z^{-2})} \in \mathfrak{R}_{s_i+s_j-3} V_k(\gf)~.
\end{equation}
Similarly, the $\mathcal{O}(z^{-1})$ pole takes the form,
\begin{equation} \label{eq:order1simp}
\begin{split}
    C^{(s_i)}(z) C^{(s_j)}(0)|_{\mathcal{O}(z^{-1})} &= \sum_{n=1}^{s_j} \sum_{m = 1}^{s_i} d_{b_1 \ldots b_{s_j}} d_{a_1 \ldots a_{m-1} c a_{m+1} \ldots a_{s_i}} if^{b_n c}_{\ \ \ \ a_m} (J^{B_{1}} J^{A} J^{B_{2}})(0) \\ &+ X^{(2)}_{a (\alpha_1 \alpha_2 \ldots \alpha_{s_i+s_j-3})} (\partial J^aJ^{\alpha_1} J^{\alpha_2} \cdots J^{\alpha_{s_i+s_j-3}})(0) + \ldots~.
\end{split}
\end{equation}
Here, we have again introduced a similar shorthand notation, 
\begin{equation}
\begin{split}
    &J^{B_{1}} = J^{b_1} J^{b_2} \cdots J^{b_{n-1}}~, \\
    &J^{A} = J^{a_1} J^{a_2} \cdots J^{a_{s_i}}~, \\
    &J^{B_{2}} = J^{b_{n+1}} J^{b_{n+2}} \cdots J^{b_{s_j}}~.
\end{split}
\end{equation}
The $\ldots$ represent operators that have at least two derivatives and are in $\mathfrak{F}_2 V_k(\gf)$. While the precise form of the second term in \eqref{eq:order1simp} can be exactly computed, it is not illuminating. The relevant key property is that $X^{(2)}_{a (\alpha_1 \alpha_2 \ldots \alpha_{s_i+s_j-3})}$ is an invariant tensor that is symmetric in the bracketed indices. As before, the symmetry follows from \eqref{eq:AntiSymJJ}, while the invariance follows from the $G$-equivariance of the OPE. 

The invariance of the tensor $d_{a_1 a_2 \ldots a_{s_i}}$ implies that $\sum_{m=1}^{s_i} f^{b c}_{\ \ a_{m}} d_{a_1 \ldots a_{m-1} c a_{m+1} \ldots a_{s_i}} = 0$, which leads to the vanishing of the first term in \eqref{eq:order1simp} above. Since it suffices to work at the level of the quotient algebra $\mathfrak{F}^1V_k(\gf)/\mathfrak{F}^2V_k(\gf)$, the problem can be recast as a question of invariant Lie theory. To this end, consider  
\begin{equation} \label{eq:F_adefn}
    F_a(j) := X_{a(\alpha_1 \alpha_2 \ldots \alpha_k)} j^{\alpha_1} j^{\alpha_2} \ldots j^{\alpha_k}~,
\end{equation}
where $X_{a(\alpha_1 \alpha_2 \ldots \alpha_k)}$ is an invariant tensor, symmetric in the bracketed indices, and $F_a(j)$ are elements in the symmetric algebra $S\gf^*$. Choosing a basis $\{e_a : a = 1,2,\ldots, \mathrm{dim}(\gf) \}$ for $\gf$ and the dual basis $\{e^a : a = 1,2,\ldots, \mathrm{dim}(\gf) \}$ for $\gf^*$, one can define $F(j) = \sum_{a=1}^{\mathrm{dim}(\gf)} F_a(j) e^a$, which is equivalently the linear map (defined on the basis of $\gf$ and extended linearly),
\begin{equation} 
\begin{split}
    F(j) &: \gf \longrightarrow S\gf^*~, \\
    &: e_a \mapsto F_a(j)~.  
\end{split}
\end{equation}
Furthermore, the invariance of the tensor $X_{a(\alpha_1 \alpha_2 \ldots \alpha_k)}$ implies that $F(j) \in \mathrm{Hom}_{\gf}(\gf, S\gf^*)$. 

By the Chevalley restriction theorem, the ring of $G$-invariant functions on $\gf$ is isomorphic to the Weyl-invariant polynomial ring on the Cartan $\mathfrak{h} \subset \gf$,
\begin{equation} \label{eq:Chevalley}
    \mathbb{C}[\gf]^G \cong \mathbb{C}[\mathfrak{h}]^W \cong \mathbb{C}[c_1,c_2,\ldots, c_{\mathrm{rank}(\gf)}]~,
\end{equation}
where $c_i$ are the independent casimir elements of $\gf$. According to Kostant's theorem \cite{Kostant:1963}, given a complex semisimple Lie algebra $\gf$, the symmetric Lie algebra $S\gf$ is a free module over the commutative ring $(S\gf)^\gf$. Moreover, for every finite dimensional representation $V$ of $\gf$, the space $\mathrm{Hom}_{\gf}(V, S \gf)$ is a free $(S\gf)^\gf$-module of rank $\mathrm{dim}(V[0])$, \emph{i.e.}, the dimension of the zero-weight space of $V$. 

The nondegenerate invariant bilinear form on $\gf$ leads to the identification $\gf \cong \gf^*$. Thus, Kostant's theorem with $V$ taken to be the adjoint representation of $\gf$, and \eqref{eq:Chevalley} together imply that $\mathrm{Hom}_{\gf}(\gf, S\gf^*)$ is a free module over $\mathbb{C}[c_1, c_2, \ldots, c_{\mathrm{rank}(\gf)}]$ of rank equal to the rank of the Lie algebra, rank$(\gf)$. Thus, $F(j)$ is expressible as,
\begin{equation} \label{eq:Kostant}
    F(j) = \sum_{i=1}^{\mathrm{rank}(\gf)} a_i f_i, \qquad \mathrm{with} \qquad a_i \in \mathbb{C}[c_1, c_2, \ldots, c_{\mathrm{rank}(\gf)}]~.
\end{equation}
In fact, a theorem attributed to Vust, identifies the basis $\{f_i\}$ with the differentials of the casimir elements, \emph{i.e.}, $f_i = dc_i$ \cite{Vust1974, PANYUSHEV2002125, panyushev2007semi}.\footnote{The author thanks Christopher Beem for helpful discussions on this topic.}

To explain this in more detail, note that the casimir elements $c_i \in (S\gf^*)^\gf$ can be equivalently regarded as polynomial maps $c_i : \gf \longrightarrow \mathbb{C}$, and the differentials of the casimir elements are then given by the linear maps $(dc_i)_x : T_{x} \gf \longrightarrow \mathbb{C}$ where $x \in \gf$. Since the tangent space $T_{x} \gf$ is isomorphic to $\gf$, the differentials are valued in $\gf^*$. Consequently, one obtains,
\begin{equation}
    (dc_i)_x = \sum_{a = 1}^{\mathrm{dim}(\gf)} \frac{\partial c_i}{\partial j^a} (x) dj^a = \sum_{a = 1}^{\mathrm{dim}(\gf)} \frac{\partial c_i}{\partial j^a} (x) e^a~. 
\end{equation}
When combined with \eqref{eq:Kostant}, this implies that $F(j)$ can be rewritten as,
\begin{equation}
    F(j) = \sum_{i=1}^{\mathrm{rank}(\gf)}\sum_{a = 1}^{\mathrm{dim}(\gf)} a_i \frac{\partial c_i}{\partial j^a} e^a,  \qquad \mathrm{with} \qquad a_i \in \mathbb{C}[c_1, c_2, \ldots, c_{\mathrm{rank}(\gf)}]~.
\end{equation}
Comparing this with the definition in \eqref{eq:F_adefn}, one finally obtains,
\begin{equation}
    X_{a(\alpha_1 \alpha_2 \ldots \alpha_k)} j^{\alpha_1} j^{\alpha_2} \ldots j^{\alpha_k} = \sum_{i=1}^{\mathrm{rank}(\gf)} \mathcal{F}_i(c_1,c_2,\ldots, c_{\mathrm{rank}(\gf)}) \frac{\partial c_i}{\partial j^a}~,
\end{equation}
where $\mathcal{F}_i$ is a polynomial function of the casimir elements. The restriction of the Casimir operators to the Zhu $C_2$ algebra is precisely given by the casimir elements as a polynomial function of the currents $J^a$. Consequently, one finds that,
\begin{equation}
    X^{(2)}_{a (\alpha_1 \ldots \alpha_{s_i+s_j-3})} (\partial J^aJ^{\alpha_1} \ldots J^{\alpha_{s_i+s_j-3}}) = \sum_{i=1}^{\mathrm{rank}(\gf)} \mathcal{F}_i\left(C^{(s_1)},\ldots, C^{(s_{\mathrm{rank}(\gf)})}\right) \partial C^{(s_i)} + \ldots~,
\end{equation}
where $\ldots$ are sums of normal ordered products of currents belonging to $\mathfrak{F}^2V_k(\gf)$, implying,
\begin{equation}
    C^{(s_i)}(z) C^{(s_j)}(0)|_{\mathcal{O}(z^{-1})} \in \mathfrak{R}_{s_i+s_j-3} V_k(\gf)~,
\end{equation}
thus establishing the goodness of the standard $\mathfrak{R}$-filtration for affine Kac--Moody algebras $V_k(\gf)$ with semisimple $\gf$. 

\subsection{Sign constraints from graded unitarity}
We now restrict our attention to affine Kac--Moody algebras $V_k(\gf)$ where $\gf$ is simply laced. Graded unitarity imposes constraints on the signs of most singular terms in the OPEs of $\mathcal{O}(z)$ and $\rho(\mathcal{O})(0)$ that are related to $\mathcal{O}$'s position in the $\mathfrak{R}$-filtration \cite{ArabiArdehali:2025fad}. In particular, writing the two-point function as
\begin{equation}
    \mathcal{O}(z) \rho(\mathcal{O})(0) = \frac{\kappa_{\mathcal{O}}}{z^{2 h_\mathcal{O}}}~,
\end{equation}
graded unitarity requires
\begin{equation}
    \kappa_\mathcal{O} \propto (-1)^{h_{\mathcal{O}} - R_{\mathcal{O}}}~,
\end{equation}
with the proportionality constant being positive. Here, we make the assumption that the $\mathfrak{R}$-filtration \eqref{eq:AKMRfiltration} is nondegenerate. Consequently, it can be refined to a unique $R$-grading by Gram-Schmidt, with the $R$-charge of the operator $\mathcal{O}(z)$ being $R_{\mathcal{O}}$. 

Determinant-based graded unitarity sign constraints for simply-laced affine Kac--Moody algebras were established in \cite{ArabiArdehali:2025fad}, and we only state the final result here. In particular, graded unitarity imposes the following sign constraint on the Gorelik--Kac vacuum determinant \eqref{eq:KacGorelikvacdet} at conformal level $L$ in a fixed weight $\nu$ sector,
\begin{equation} \label{eq:grunitsignrule}
    \mathrm{sign}\left(\mathrm{det}S_{L\delta+\nu} \right) = (-1)^{\sum_{\substack{p \geqslant 2, q \geqslant 1:} (p,q) = 1} m_{p/q}(L\delta + \nu)}~.
\end{equation}

For any interacting $\mfR$-filtration of four-dimensional type, $R(J) = R(T) = 1$. Since the stress tensor for $V_k(\gf)$ is given by the Segal--Sugawara stress tensor, graded unitarity requires $-h^\vee < k \leqslant 0$, and thus, the factors of $(k+h^\vee)^{m_0}$ in the vacuum determinant are positive. Hence, \eqref{eq:grunitsignrule} implies that when the vacuum determinant in a fixed-weight sector is non-vanishing, then only an \emph{even} number of non-critical factors $(k - k_{p,q})$ are positive.

Before implementing graded unitarity constraints for simply-laced affine Kac--Moody algebras, we briefly outline the general strategy. At a given conformal level $L$, we use the closed-form vacuum determinant \eqref{eq:KacGorelikvacdet} at a fixed weight $\nu$ to determine the \emph{structure of zeros}. More precisely, this involves determining the most negative and second most negative zeros (in certain cases, also the third most negative zero) of the vacuum determinant, along with their multiplicities. While we primarily consider the ``fixed-root sector" \emph{i.e.}, $\nu = \alpha$, where $\alpha$ is any root of the respective Lie algebra $\gf$, other weight sectors are also analogously examined at certain low levels. Determinant-based graded unitarity constraints from \eqref{eq:grunitsignrule} for a given conformal level $L$ and an appropriate fixed weight $\nu$ are then imposed to exclude intervals of levels $k$ incompatible with graded unitarity. In addition, a bespoke low-level analysis by passing to a representation basis instead of a weight basis facilitates the exclusion of certain isolated levels that survive the examined determinant-based constraints. Extending this procedure to all conformal levels $L$, we find that the closed-form vacuum determinant \eqref{eq:KacGorelikvacdet}, together with the graded unitarity sign constraint \eqref{eq:grunitsignrule}, provides a remarkably powerful criterion for excluding levels $k$ of all simply-laced affine Kac--Moody algebras, ultimately leading to the classification summarised in Table \ref{tab:classification}.

